\subsection{The inverse}
\label{icg:sec:inverse}

\subparagraph{Block structure and the projector.}
Write $\mc H^\downarrow := \ker b_0^+ \cap \ker b_0^-$ and $W := c_0^+ c_0^- \mc H^\downarrow$, with $c_0^\pm = \tfrac12 (c_0 \pm \bar c_0)$, so that $\mc H = \mc H^\downarrow \oplus c_0^+ \mc H^\downarrow \oplus c_0^- \mc H^\downarrow \oplus W$.
In the eigenbasis $\ket{h, m}$ of $(L_0^+, L_0^-)$ the operator of~\eqref{p:eq:ws:p-circ} is
\begin{gather}
  \label{icg:eq:G}
  G
    := - \I \kappa \, b_0^+ b_0^- \, F,
  \qquad
  F \ket{h, m}
    = \hat f_m(h) \ket{h, m},
  \\
  \label{icg:eq:fhat}
  \hat f_m(h)
    := \hat f_{Q_V}(h, m),
  \qquad
  \hat f_R(h, m)
    := \int_{R} \dd s \, \dd\theta \; \e^{-s h + \I m \theta},
\end{gather}
the Laplace transform in the neck length and Fourier transform in the twist of the region.
$G$ maps $W$ to $\mc H^\downarrow$ and annihilates the other three summands, so the form $\braket{\cdot, G \, \cdot}$ lives on $W \times W$, non-degenerate wherever $\hat f_m(h) \neq 0$; as a form on $\mc H^{\wedge 2}$, $p = \Phi_2^\dagger G \Phi_2$ with $\Phi_2^\dagger$ the transpose of the circle map for the pairing.

\subparagraph{The projector.}
The inverse $p^{-1}$ is defined on $\mc H^{\wedge 2} / \ker p$; pushing it back to $\mc H$ requires a choice of section, and this is where the kernel of $p$ enters.
With such a choice, $\Pi := \Phi_2\, p^{-1} \Phi_2^\dagger$.
From $p = \Phi_2^\dagger G \Phi_2$ one line gives $G \Pi G = G$ on the image $\mc I$ of the circle states, when paired against $\mc I$.
Moreover, $P_{\mc I} := \Pi G$ is the $G$-orthogonal projector onto $\mc I$, modulo $\mc I \cap \mc I^{\perp_G}$.
Since $G$ sees only the $W$-components of the circle states,
\begin{equation}
  \label{icg:eq:Pi}
  \Pi
    = P_{\mc I} \, G_W^{-1},
\end{equation}
with $G_W^{-1}$ the inverse of $G$ on the block, $\mc H^\downarrow \to W$: the inverse of the restricted form is the block inverse followed by a projection, and the projection is the whole difference.
It is invisible when both neighbors are circle states (\cref{p:sec:ws:telescoping}).

\subparagraph{The hypothesis.}
$G P_{\mc I} = G$ on all of $W$ – the projector is invisible to the gluing – only under a hypothesis: that the projection of the circle states onto $W$, extended to $\Sigma_2$ of unbounded level (the space in which the field lives), reaches the whole even-$L_0^-$, ghost-number-4 part of $W$ at fixed total momentum, modulo $\ker G$.
It would be falsified by a persistent rank deficit of the projected circle map, level by level, against the dimension of the even block.
This is motivated – at fixed $P$ the inputs include arbitrary functions of the relative momentum and all oscillator states at both punctures while each level subspace of $W$ is finite-dimensional, and the matter part of the map is triangular in level with an invertible leading part – but the ghost sector is the expected failure mode, since $c(x) c(-x)$ starts at $c \pd c$.
Under the hypothesis $p^{-1}$ acts on $\Sigma_2$ of unbounded level, outside the finite-level setting in which every pairing of circle states is a finite correlator; density of the finite-level image would not suffice, because $\Phi_2^\dagger$ sums over all levels.

\subparagraph{The two $c$'s are the circle's zero modes.}
On $\mc H^\downarrow$, with $\{b_0^\pm, c_0^\pm\} = 1$, $\{b_0^\pm, c_0^\mp\} = 0$ and $b_0^\pm \ket{\downarrow} = 0$,
\begin{equation}
  \label{icg:eq:zero-modes}
  b_0^+ b_0^- \, c_0^- c_0^+ \ket{\downarrow}
    = b_0^+ (1 - c_0^- b_0^-) c_0^+ \ket{\downarrow}
    = \ket{\downarrow},
  \qquad
  c_0^- c_0^+ \, b_0^+ b_0^- \big|_W
    = 1,
\end{equation}
so $(b_0^+ b_0^-)^{-1} = c_0^- c_0^+$ between $W$ and $\mc H^\downarrow$, and in the fixed-halves model
\begin{equation}
  \label{icg:eq:inverse}
  \Phi_2 \, (p^{\rm circ}_{2,2})^{-1} \, \Phi_2^\dagger
    = P_{\mc I} \; \frac{\I}{\kappa} \; c_0^- c_0^+ \; F^{-1},
  \qquad
  F^{-1} \ket{h, m}
    = \frac{1}{\hat f_m(h)} \ket{h, m}.
\end{equation}
The two $c$'s of $p^{-1}$ are the zero modes dual to the two Schiffer antighosts of the gluing, i.e. to the dilation and the rotation of the circle: the residual $\C^*$ of the two-punctured sphere, the pair of conformal Killing directions the counting asks for.
This identification is a fixed-halves statement: the dressed antighosts~\eqref{p:eq:ws:dressing} contain $b_n$ with $n \neq 0$ in the circle frame, and for $p_{2,2}$ itself only the ghost-number count transfers.

\subparagraph{Not a gluing over a reversed region.}
The group law $q^{L_0} \bar q^{\bar L_0} \, q'^{L_0} \bar q'^{\bar L_0} = (qq')^{L_0} (\bar q \bar q')^{\bar L_0}$ holds pointwise, but the $b$-dressed gluings $\mc G_R = \int_R \dd s \, \dd\theta \, b_0^+ b_0^- \e^{-s L_0^+} \e^{\I\theta L_0^-}$ compose to $\mc G_R \mc G_{R'} \propto (b_0^+ b_0^-)^2 = 0$: a gluing absorbs two units of ghost number through its two antighosts, an inverse must supply them, so it carries $c$'s rather than a reversed parameter.
At the level of the modulus, the inverse of an integral over a region is not the integral over the inverse region: for one sector and a constant cutoff $\sigma$,
\begin{equation}
  \label{icg:eq:no-inverse-region}
  \hat f_{[-\sigma, 0]}(h) \; \hat f_{[0, \sigma]}(h)
    = \frac{(\e^{\sigma h} - 1)(1 - \e^{-\sigma h})}{h^2}
    \neq 1 .
\end{equation}
Only a point inverts to a point (a frozen neck to the frozen neck of opposite length), and a two-form has no point to invert.

\subparagraph{The reciprocal of the region's transform.}
For a rotation-independent cutoff, $Q = \{1 < \abs{q} < \e^{\sigma}\}$, the twist integral gives $2\pi \delta_{m,0}$ and the neck integral the transform of the slab $[-\sigma, 0]$,
\begin{equation}
  \label{icg:eq:tower}
  \hat f_m(h)
    = 2\pi \, \delta_{m,0} \, \frac{\e^{\sigma h} - 1}{h},
  \qquad
  \frac{1}{\hat f_0(h)}
    = \frac{1}{2\pi} \, \frac{h}{\e^{\sigma h} - 1}
    = \frac{h}{2\pi} \sum_{n \ge 1} \e^{-n \sigma h}
  \quad (\Re h > 0),
\end{equation}
the other branch, $\Re h < 0$, being the one we call reversed.
Since $h \, \e^{-n\sigma h} = \int \dd s \, \delta'(s - n\sigma) \, \e^{-sh}$, the reciprocal is the Laplace transform of $\sum_{n \ge 1} \delta'(s - n\sigma)$: a discrete tower of necks of positive lengths $n\sigma$, each carrying one derivative in its length.
In operator form each term is $c_0^- c_0^+ L_0^+ \e^{-n\sigma L_0^+}$ with $L_0^+ = \{Q_B, b_0^+\}$, the zero-mode part of the $\Psi$ kinetic operator $c_0^- Q_B$ on a cylinder of length $n\sigma$; the $L_0^+$ is $-\pd_s$ of the plumbing family, a derivative in the neck length acting on whatever is glued to the cylinder.

Up to $c_0^- c_0^+$ in place of $Q_B$ this is the kinetic operator $Q_B L_0^+ / (\e^{\sigma L_0^+} - 1)$ of an auxiliary field with one string leg, obtained by gluing over the reversed plumbing range: between the $\Sigma_2$ formulation and the single-string one the roles of region and tower are exchanged, the tower being the kinetic operator there and the region its propagator.
The same exchange appears for the string field itself: the zero-mode part of the $\Psi$ kinetic form is $c_0^- c_0^+ L_0^+$, the point $s = 0$ with a $\delta'$ (transform $h$), and its inverse is the half-line $[0, \infty)$ (transform $1/h$); the half-line inverts to a single $\delta'$ because it is closed under addition, while a bounded slab inverts to the $\delta'$ at one endpoint convolved with the geometric series generated by the other – that series is the whole difference between $Q_B^{-1}$ and $(p^{\rm circ}_{2,2})^{-1}$ at constant cutoff.

In the level-matched sector the reciprocal~\eqref{icg:eq:tower} is regular at $h = 0$ (value $1/(2\pi\sigma)$): the $\Sigma_2$ line has no mass-shell pole, as for the open-string auxiliary of~\cite{ErbinFirat2023}.
Its poles sit at $h = 2\pi \I k / \sigma$, $k \neq 0$, the zeros of $\hat f_0$: complex momenta, zeros of the block operator $G_W$ on that sector, and poles of $p^{-1}$ only under the hypothesis above (one sector vanishing does not make $\det p$ vanish, every level of a circle state being populated).

The actual cutoff $\sigma_V(\theta)$ is not constant; then $\hat f_m(h)$ is the $m$-th Fourier coefficient of $(\e^{\sigma_V(\theta) h} - 1)/h$, and odd $m$ vanish by $\pi$-periodicity, consistently with the circle states having no odd-$L_0^-$ component.
In the level-matched sector, with $\nu$ the distribution of $\sigma_V(\theta)$ for $\theta$ uniform on $[0, 2\pi)$, a probability measure on $[\ln\tfrac94, \ln 9]$ at $\beta = 3/2$, and $\hat\nu(h) := \int \dd\nu(\sigma) \, \e^{\sigma h}$,
\begin{equation}
  \label{icg:eq:tower-sigmaV}
  \frac{1}{\hat f_0(h)}
    = \frac{1}{2\pi} \, \frac{h}{\hat\nu(h) - 1}
    = - \frac{h}{2\pi} \sum_{n \ge 0} \hat\nu(h)^n
  \quad (\Re h < 0),
\end{equation}
a tower of cylinders whose lengths are sums of $n$ draws from the seam distribution, on the reversed branch only (on the other branch there is a tower only if $\nu$ is a point mass); none of its zeros is real, since $(\hat\nu(h) - 1)/h > 0$ for real $h$.
In the even sectors $m \neq 0$, however, $\hat f_m(h) = \frac1h \int_0^{2\pi} \e^{\I m\theta} \e^{\sigma_V(\theta) h} \, \dd\theta$ is real for real $h$ (the profile is even in $\theta$), and its asymptotic signs are fixed by the geometry: as $h \to +\infty$ the flat fold arcs $(\theta_1, \theta_2)$ dominate, with sign $\operatorname{sgn}[\sin m\theta_2 - \sin m\theta_1]$, positive for $m = 4$, while as $h \to -\infty$ the minima at $\theta = 0, \pi$ dominate and every even $\hat f_m$ is negative.
So $\hat f_4$ has an odd number of real zeros.
Numerically (at $\beta = 3/2$, from the computed profile, three digits, scanning $h \in [-10, 10]$), $\hat f_4$ has one, at $h_* \approx -2.36$, $\hat f_6$ has two, at $-4.89$ and $-1.91$, and $\hat f_2$ has none.
Since $h$ is the $L_0^+$ eigenvalue of the circle state, $h = \tfrac{\alpha'}{2} P^2 + N + \bar N - 2$ with $P$ the momentum through the circle, these are zeros of $G_W$ at \emph{real timelike} momentum through the circle, $\tfrac{\alpha'}{2} (-P^2) = N + \bar N - 2 - h_*$ (at six points, in the channel invariant $s_{123}$ of \cref{icg:sec:six-points}), in sectors $N - \bar N = m \neq 0$ that contain no level-matched state; they are never reached at real Euclidean momentum, where $h \ge N + \bar N - 2$, which is $\ge 2$ in the $m = 4$ sector.
Under the hypothesis they are poles of $p^{-1}$; their fate in the composite is open (\cref{icg:sec:six-points}).


